\chapter{Základné pojmy}
\label{chap:terms}

Vzhľadom na nekonzistentnosť pojmov naprieč štandardami, literatúrou a legislatívou by sme chceli zjednotiť tie pojmy,
ktoré budeme v nasledujúcich kapitolách používať, aby sme si s čitateľom spolu porozumeli. V prípade, že čitateľ je
znalý v problematike, môže preskočiť priamo na ďalšiu kapitolu \ref{chap:it-grundschutz} -- \nameref{chap:it-grundschutz},
pričom sa môže vrátiť k tejto kapitole, keby mal nejaké nejasnosti. Vychádzať budeme zo slovníka \cite{KUDIKB-MVS} a niektoré
pojmy preberieme priamo z \cite{BP}.

\begin{itemize}
  \item \textbf{Kybernetická a informačná bezpečnosť (KIB)}, multidisciplinárna oblasť ľudskej činnosti, zaoberajúca sa skúmaním
    hrozieb voči informáciám, informačným systémom a sieťam a návrhom riešení, ktoré by naplneniu hrozieb zamedzili,
    alebo aspoň zmiernili ich dopad. V tejto práci sú pre nás pojmy kybernetická bezpečnosť, informačná bezpečnosť
    a kybernetická a informačná bezpečnosť zameniteľné a znamenajú to isté.

  \item \textbf{Aktívum}, angl. \textit{asset}, je čokoľvek, čo má pre organizáciu hodnotu, môžu byť cieľom útoku a má zmysel 
    (resp. je potrebné) chrániť pred potenciálnymi hrozbami. Aktíva sú hmotné (zariadenia, infraštruktúra, personál) a nehmotné
    (peniaze, informácie, vedomosti).

  \item \textbf{Hrozba}, angl. \textit{threat}, je objektívne existujúca potenciálna možnosť priamo alebo nepriamo narušiť
    (alebo negatívne ovplyvniť) systém, spracovávané informácie, alebo iné aktíva organizácie.

  \item \textbf{Dôvernosť}, angl. \textit{confidentiality}, je bezpečnostná požiadavka, ktorej naplnenie znamená, že informáciu
    obsiahnutú v správe sa nedozvedia nepovolené osoby.

  \item \textbf{Integrita}, angl. \textit{integrity}, je bezpečnostná požiadavka, ktorej naplnenie znamená, že údaje nie je
    možne zmeniť bez toho, aby to ich vlastník alebo adresát nevedel zistiť.

  \item \textbf{Dostupnosť}, angl. \textit{availability}, je bezpečnostná požiadavka, ktorej naplnenie znamená, že zdroje systému
    sú k dispozícií oprávnenej - osobe v momente, keď o to požiada; alebo od istého okamihu; alebo v istom časovom rámci
    (napr. dni v týždni, časový interval počas dňa).

  \item \textbf{CIA}, skratka, ktorá zlučuje trojicu dôvernosť, integrita a dostupnosť, sú to základné hodnoty aktíva, ktoré
    sa usilujeme chrániť.
  
  \item \textbf{Riziko}, angl. \textit{risk}, je veličina závislá od závažnosti (možného dopadu) hrozby a pravdepodobnosti, 
    že sa hrozba naplní.

  \item \textbf{Analýza rizík}, angl. \textit{risk analysis}, proces identifikácie rizík, ohodnotenia rizík a stanovenia
    úrovne rizík.
    
  \item \textbf{Bezpečnostné opatrenie}, prostriedok, ktorým znižuje pravdepodobnosť alebo potenciálny dopad hrozby,
  čím znížime jej výsledné riziko.

  \item \textbf{Vedenie organizácie}, skrátene \textit{vedenie}, samostatná osoba alebo skupina ľudí, ktorá je zodpovedná 
  za riadenie a kontrolu organizácie na najvyššej úrovni (napr. riaditeľ, generálny riaditeľ, predstavenstvo, štatutárny orgán).

  \item \textbf{Manažér kybernetickej a informačnej bezpečnosti}, skrátene \textit{m-KIB}, organizačná rola, osoba, ktorá je zodpovedná
    za manažment a koordináciu bezpečnostného procesu v organizácii.

\end{itemize}